\documentclass[11pt,a4paper]{article}

\usepackage[utf8]{inputenc}
\usepackage[T1]{fontenc}
\usepackage[margin=2.5cm]{geometry}
\usepackage{amsmath,amssymb,amsthm}
\usepackage{booktabs,array,longtable,tabularx,enumitem}
\usepackage{microtype}
\usepackage{xcolor}
\usepackage[colorlinks=true,
            linkcolor=blue!55!black,
            citecolor=blue!55!black,
            urlcolor=blue!55!black]{hyperref}
\usepackage{url}

\hypersetup{
  pdftitle={HypG: A Safe-Rust Post-Quantum Cryptographic Suite},
  pdfauthor={HypG Project},
  pdfsubject={White Paper: design, architecture, and evaluation of the v2.0 research profile},
  pdfkeywords={post-quantum cryptography, Module-LWE, BLAKE3, ChaCha20, FORS, WOTS+, XMSS, safe Rust}
}

\setlength{\emergencystretch}{2.5em}
\allowdisplaybreaks

\theoremstyle{plain}
\newtheorem{theorem}{Theorem}[section]
\newtheorem{proposition}[theorem]{Proposition}
\newtheorem{lemma}[theorem]{Lemma}
\newtheorem{corollary}[theorem]{Corollary}
\theoremstyle{definition}
\newtheorem{definition}[theorem]{Definition}
\newtheorem{assumption}[theorem]{Assumption}
\theoremstyle{remark}
\newtheorem{remark}[theorem]{Remark}

\newcommand{\concat}{\,\Vert\,}
\newcommand{\bits}{\{0,1\}}
\newcommand{\Z}{\mathbb{Z}}
\newcommand{\Ring}{\mathcal{R}_q}
\newcommand{\Hash}{\mathsf{H}}
\newcommand{\XOF}{\mathsf{XOF}}
\newcommand{\KDF}{\mathsf{KDF}}
\newcommand{\PRF}{\mathsf{PRF}}
\newcommand{\KeyGen}{\mathsf{KeyGen}}
\newcommand{\Encaps}{\mathsf{Encaps}}
\newcommand{\Decaps}{\mathsf{Decaps}}
\newcommand{\Encode}{\mathsf{Encode}}
\newcommand{\Decode}{\mathsf{Decode}}
\newcommand{\Pack}{\mathsf{Pack}}
\newcommand{\Unpack}{\mathsf{Unpack}}
\newcommand{\Sign}{\mathsf{Sign}}
\newcommand{\Verify}{\mathsf{Verify}}
\newcommand{\negl}{\mathsf{negl}}
\newcommand{\poly}{\mathsf{poly}}
\newcommand{\ppt}{\mathsf{PPT}}
\newcommand{\adv}{\mathcal{A}}
\newcommand{\chall}{\mathcal{C}}
\newcommand{\ro}{\mathcal{O}}
\newcommand{\concatx}{\concat}

\title{\textbf{HypG: A Safe-Rust Post-Quantum Cryptographic Suite}\\[0.5ex]
\large White Paper --- Design, Architecture, and Evaluation of the v2.0 Research Profile\\[0.3ex]
\normalsize (KEM v2.0; Symm/Hash v2.0; DS v2.0.1)}

\author{HypG Project\\
\small Independent research prototype}
\date{October 2026}

\begin{document}
\maketitle

\begin{abstract}
HypG is a from-scratch, zero-external-dependency cryptographic suite
written in Rust with \texttt{\#![forbid(unsafe\_code)]} in every
library.  It comprises a Module-LWE key-encapsulation mechanism, a
portable BLAKE3 hash/XOF/keyed-hash layer, a ChaCha20 plus keyed-BLAKE3
authenticated-encryption scheme (one-shot and streaming), and a
stateless hash-based signature scheme.  The v2.0 profile uses module
rank four over $\Z_{4096}[X]/(X^{256}+1)$ for the KEM and an
SLH-DSA-256f-shaped hypertree with FORS $35\times 9$ for signatures.
The signature wire format carries the separate \textbf{2.0.1}
revision: explicit PRF/F domain tags and FIPS-205-style odd-node
promotion in the WOTS+ L-tree.

This paper documents the design rationale, formal construction,
implementation discipline, controlled evaluation, and---most
importantly---the limits of the evidence.  The recorded release has 146
deterministic tests passing in debug, release, and on the Rust 1.70
MSRV; four frozen KAT files totalling 234,114 bytes; a 14-suite
adversarial harness; a 20,000-trial failure sweep; a full 11,520-bit
sweep of one test ciphertext; Miri coverage of the KEM, AEAD, and DS bounds
tests; and byte-identical KAT output on x86-64 and
\texttt{qemu-aarch64}.  None of this is a security proof.  The KEM's
exact FO-style transform is not covered by a published theorem, the
lattice hardness of the chosen parameters has not been confirmed by an
external estimator, the signature composition has no exact
signing-volume bound, and there is no external audit, formal
verification, \texttt{dudect}-grade timing study, or power/EM/fault
campaign.  HypG v2.0 is a research prototype and must not be deployed
to protect real assets.
\end{abstract}

\noindent\textbf{Keywords:} post-quantum cryptography; Module-LWE;
Fujisaki--Okamoto; hash-based signatures; FORS; WOTS+; XMSS; BLAKE3;
ChaCha20; authenticated encryption; safe Rust; reproducible evaluation.

\medskip
\noindent\textbf{Document status.} Technical report accompanying the
v2.0 artifact.  It is not peer-reviewed, not an audit, and not a
security certification.

\tableofcontents
\newpage

%======================================================================
\section{Introduction}
\label{sec:intro}

Public-key cryptography deployed today is vulnerable to a
cryptographically relevant quantum computer: Shor's algorithm solves
integer factorisation and discrete logarithms in polynomial time.
Symmetric primitives are less affected, but their effective security
halves under Grover search, so a post-quantum suite must combine
quantum-resistant public-key constructions with conservative symmetric
cryptography and a disciplined implementation.

HypG is an experiment in building the \emph{whole} suite rather than a
single primitive: key establishment, hashing, authenticated encryption,
and signatures, sharing one small implementation core.  The project
imposes unusually strict engineering constraints: no external
dependencies, no \texttt{unsafe} code, \texttt{no\_std} cryptographic
libraries, fixed public loop bounds, explicit little-endian encodings,
deterministic test vectors, and mechanically checked CI gates.

The current release is \textbf{v2.0.0} with a \textbf{DS v2.0.1} wire
revision.  It is a research prototype.  The documentation deliberately
separates (i) what the code does, (ii) what the test suite establishes,
and (iii) what would still be required for a security claim.  This
paper follows the same separation.

\subsection{Contributions}

The contributions are:

\begin{enumerate}[leftmargin=*]
\item \textbf{A complete, dependency-free implementation.}  A
  \texttt{no\_std}, no-allocation cryptographic core and four protocol
  pillars, with all dependency edges confined to workspace path
  dependencies and all library crates under
  \texttt{\#![forbid(unsafe\_code)]}.

\item \textbf{A precise v2.0 parameter profile.}  A rank-four
  Module-LWE KEM at lattice dimension 1024 with a power-of-two modulus
  $q=4096$, and an SLH-DSA-256f-shaped signature profile with
  $n=32$, $h=68$, $d=17$, $h'=4$, FORS $(35,9)$, and WOTS+ $w=16$.

\item \textbf{A documented DS 2.0.1 revision.}  Explicit one-byte
  $\PRF$/$F$ tags remove a structural collision between the two
  64-byte keyed-hash inputs, and odd WOTS+ L-tree nodes are promoted
  unchanged in the FIPS 205 style.

\item \textbf{An honest treatment of the CA-PRF.}  The cellular
  automaton is presented as a branchless mixing and domain-separation
  stage whose uniformity argument rests on the final BLAKE3 XOF, not
  as an independently proven hardness assumption.

\item \textbf{A corrected KEM correctness analysis.}  The compression
  residuals from $c_1$ and $c_2$ are included explicitly, yielding a
  working standard deviation of approximately 70 and a margin of
  approximately $14.6\sigma$ against the 1024 decision boundary.  This
  is a heuristic estimate, not a formal failure bound.

\item \textbf{A byte-identical streaming AEAD.}  The streaming front
  uses the same subkeys, MAC input, counter schedule, and tag as the
  one-shot construction; only the scheduling differs.

\item \textbf{A reproducible evaluation record.}  Frozen KAT checksums,
  cross-architecture equality, a deterministic red-team harness, a
  local lattice lab, Miri subsets, and a controlled release benchmark
  are recorded with commands and machine details.

\item \textbf{Explicit non-claims and an audit agenda.}  The paper names
  the exact proof, estimator, side-channel, and operational gaps that
  stand between the current artifact and a deployable system.
\end{enumerate}

\subsection{What this paper does not claim}

This paper does \emph{not} claim that HypG is IND-CCA2 secure,
UF-CMA/SUF-CMA secure, or a secure AEAD in the sense of a published
reduction for its exact composition.  It does not claim NIST Category-5
security, FIPS compliance, certification, or production readiness.  It
does not claim constant-time execution at the machine-code level, and
it does not claim resistance to power, electromagnetic, or fault
attacks.  KAT vectors pin behaviour; they are not security evidence.

\subsection{Organisation}

Section~\ref{sec:background} reviews the relevant background.
Section~\ref{sec:principles} states the design constraints.
Section~\ref{sec:architecture} describes the architecture.
Sections~\ref{sec:kem}--\ref{sec:ds} describe the four pillars.
Section~\ref{sec:implementation} discusses implementation assurance.
Section~\ref{sec:evaluation} records the controlled evaluation.
Section~\ref{sec:security} analyses assumptions and open problems.
Section~\ref{sec:validity} discusses threats to validity.
Section~\ref{sec:related} relates HypG to prior work.
Section~\ref{sec:roadmap} gives the path to a defensible profile.

%======================================================================
\section{Background and related work}
\label{sec:background}

\subsection{Post-quantum cryptography}

The NIST post-quantum standardisation process produced ML-KEM
(FIPS~203)~\cite{fips203}, ML-DSA (FIPS~204)~\cite{fips204}, and
SLH-DSA (FIPS~205)~\cite{fips205}.  ML-KEM is a Module-LWE KEM;
ML-DSA is a Module-LWE/Module-SIS signature scheme; SLH-DSA is a
stateless hash-based signature scheme derived from
SPHINCS+~\cite{sphincs}.  HypG does not implement any of these
standards verbatim.  It uses the same broad mathematical families and
some of the same dimension choices, but with different primitives,
encodings, domain separation, and wire formats.

\subsection{Module-LWE}

Module-LWE generalises ring-LWE by working over $R_q^k$.  For a public
matrix $A$ and secret/error vectors $s,e$ drawn from a small
distribution, the public key is $b=A s+e$.  Encryption uses a fresh
small vector $r$ and errors $e_1,e_2$:
\[
  c_1 = A^{T}r+e_1, \qquad
  c_2 = \langle b,r\rangle + e_2 + \Encode(m).
\]
The decryption noise is $\langle e,r\rangle+e_2-\langle s,e_1\rangle$;
correctness requires it to remain inside the decoding region.  The
concrete hardness depends on the ring degree, module rank, modulus,
noise distribution, and number of samples.  The community
lattice-estimator~\cite{estimator} is the standard tool for estimating
the cost of known attacks; it has not been run on HypG's exact
parameters.

\subsection{Hash-based signatures}

SLH-DSA and SPHINCS+ combine a few-time FORS signature with a hypertree
of XMSS/WOTS+ one-time signatures.  The hash function is tweaked by a
public seed and an address so that every hash call is
domain-separated.  HypG follows this architecture with BLAKE3 as the
tweakable hash, a custom 32-byte address, and a deterministic
randomizer.  Determinism is useful for reproducible vectors but removes
fresh-randomness hedging; the security consequences are discussed in
Section~\ref{sec:ds}.

\subsection{Symmetric primitives}

BLAKE3~\cite{blake3} is a tree-structured hash and XOF with a
32-byte keyed mode and a reference \texttt{derive\_key} mode.
ChaCha20~\cite{rfc8439} is a 20-round ARX stream cipher with a 96-bit
nonce and a 32-bit block counter.  HypG uses both as designed: BLAKE3
for hashing, expansion, key derivation, and MACs; ChaCha20 for
confidentiality.  The AEAD is an Encrypt-then-MAC composition rather
than a standardised AEAD algorithm.

%======================================================================
\section{Design principles}
\label{sec:principles}

Five constraints shape every component.

\begin{enumerate}[leftmargin=*]
\item \textbf{No external dependencies.}  Every dependency edge is a
  workspace path dependency.  The supply-chain surface is therefore the
  Rust toolchain and the repository itself, not a transitive crate
  graph.  \texttt{cargo tree} is part of the release gate.

\item \textbf{No \texttt{unsafe}.}  Every library crate uses
  \texttt{\#![forbid(unsafe\_code)]}.  Arithmetic and conversion paths
  are written to be total where possible.  The cost is scalar code and
  best-effort, rather than volatile, zeroization.

\item \textbf{\texttt{no\_std} cryptographic libraries.}  The core,
  KEM, Symm, AEAD, DS, and hash crates do not require an allocator, a
  filesystem, or an operating system.  Threads, files, and allocation
  live only in \texttt{std} crates.

\item \textbf{Fixed public geometry.}  Protocol loops have public
  bounds.  Secret-dependent branches, secret-dependent array indices,
  and secret-dependent early returns are avoided in the cryptographic
  paths.

\item \textbf{Reproducible evidence.}  Deterministic KATs, a frozen
  checksum inventory, an adversarial harness, a lattice lab,
  cross-architecture equality, and a recorded benchmark are part of the
  artifact.
\end{enumerate}

These are engineering constraints, not security proofs.  In particular,
source-level branchlessness does not imply machine-code constant-time
execution, and safe Rust does not imply resistance to microarchitectural
leakage.

%======================================================================
\section{Architecture}
\label{sec:architecture}

\subsection{Component graph}

The dependency graph is acyclic:
\[
\begin{array}{c}
\texttt{hypg-core}\\[-1mm]
\swarrow\ \downarrow\ \searrow\\[-1mm]
\texttt{hypg-kem}\quad
\texttt{hypg-symm}\to\texttt{hypg-aead}\quad
\texttt{hypg-ds}\to\texttt{hypg-mt}
\end{array}
\]
\texttt{hypg-core} supplies portable BLAKE3, ChaCha20, the CA-PRF
expander, polynomial arithmetic over $\Ring$, and constant-time
helpers.  \texttt{hypg-hash} is a thin API over the core hash
functions.  \texttt{hypg-kat}, \texttt{hypg-bench},
\texttt{hypg-redteam}, \texttt{hypg-lattice}, and \texttt{hypg-cli} are
\texttt{std} harnesses or operator tools.

\subsection{Version map}

\begin{table}[t]
\centering
\small
\begin{tabular}{@{}llll@{}}
\toprule
Component & Code version & Wire/spec revision & Frozen vectors \\
\midrule
Workspace & 2.0.0 & --- & --- \\
\texttt{hypg-kem} & 2.0.0 & KEM v2.0 & \texttt{kats/v2.0.0.txt} \\
\texttt{hypg-symm} & 2.0.0 & Symm v2.0 & \texttt{kats-symm/v2.0.0.txt} \\
\texttt{hypg-hash} & 2.0.0 & Hash v2.0 & \texttt{kats-hash/v2.0.0.txt} \\
\texttt{hypg-ds} & 2.0.0 & \textbf{DS v2.0.1} & \texttt{kats-ds/v2.0.1.txt} \\
\bottomrule
\end{tabular}
\caption{Version identity of the active artifacts.  The literal
\texttt{HYP-G/v07/...} strings are frozen domain separators, not a
statement that the parameter set is v0.7.}
\label{tab:versions}
\end{table}

\subsection{Boundary decisions}

Threading is deliberately excluded from the \texttt{no\_std} crates.
\texttt{hypg-mt} uses \texttt{std::thread::scope} and owns all worker
results; deterministic assembly makes its output byte-identical to the
serial signer.  SIMD intrinsics are excluded by the \texttt{unsafe}
policy and portable SIMD is not stable on the MSRV; the implementation
is scalar and relies on the compiler only for ordinary code generation.

%======================================================================
\section{The Module-LWE KEM}
\label{sec:kem}

\subsection{Ring and parameters}

Let
\[
  \Ring=\Z_q[X]/(X^{256}+1), \qquad q=2^{12}=4096,
\]
and let $k=4$.  Secrets and errors are drawn from $\mathrm{CBD}(2)$,
whose coefficients lie in $\{-2,-1,0,1,2\}$ with probabilities
$1/16,4/16,6/16,4/16,1/16$ and variance one.  The message is 32 bytes,
encoded one bit per coefficient at spacing $q/2=2048$.  The public key
is
\[
 pk=\mathit{seed}_A(32)\concat \Pack_{12}(b_0)\concat\cdots\concat
 \Pack_{12}(b_3),
\]
the ciphertext is
\[
 ct=\Pack_{10}(\mathrm{Compress}_{10}(c_{1,0}))\concat\cdots\concat
 \Pack_{10}(\mathrm{Compress}_{10}(c_{1,3}))
 \concat\Pack_5(\mathrm{Compress}_5(c_2)),
\]
and the secret key is $sk=\mathit{seed}\concat\mathit{fake}\concat pk$.
The byte sizes are $1568$, $1440$, $1632$, and $32$ for
$pk,ct,sk,ss$.

\subsection{Key generation and encapsulation}

Key generation derives $\mathit{seed}_A$ from the master seed with the
CA-PRF, expands a $4\times4$ matrix $A$ from it, expands $s$ and $e$
from the master seed, and computes $b=A s+e$.  A separate CA-PRF
domain derives the 32-byte implicit-rejection key $\mathit{fake}$.

Encapsulation treats the caller-provided 32-byte \texttt{coins}
argument as the FO message $m$:
\[
 \mathit{rseed}=\Hash(m\concat pk\concat\texttt{"HYP-G/v07/r-seed"}),
\]
then derives $(r,e_1,e_2)$ from $\mathit{rseed}$ using the
\texttt{rand-r}, \texttt{rand-e1}, and \texttt{rand-e2} domains.  It
computes
\[
 c_1=A^T r+e_1,\qquad
 c_2=\langle b,r\rangle+e_2+\Encode(m),
\]
packs the ciphertext, and derives
\[
 ss=\XOF(\texttt{"HYP-G/v07/ss"}\concat m\concat ct\concat
 \Hash(pk),32).
\]

\paragraph{API warning.}  The \texttt{coins} value is not an auxiliary
nonce.  It is the message itself and the input to the FO randomness
generator.  Reusing it under the same public key repeats both the
ciphertext and the shared secret.  It must be a fresh CSPRNG output for
every encapsulation.

\subsection{Decapsulation and implicit rejection}

Decapsulation unpacks and decompresses $(c_1,c_2)$, recomputes
$s$ from the secret seed, and obtains
\[
 v=c_2-\langle s,c_1\rangle,\qquad m'=\Decode(v).
\]
It then recomputes the encapsulation randomness and the full ciphertext
$ct'$, compares $ct'$ with the received $ct$ using a full-scan
constant-time equality, and computes
\[
 ss_{\mathrm{ok}}=\XOF(\texttt{".../ss"}\concat m'\concat ct\concat
 \Hash(pk),32), \qquad
 ss_{\mathrm{rej}}=\XOF(\texttt{".../reject"}\concat\mathit{fake}\concat
 ct,32).
\]
A branchless select returns $ss_{\mathrm{ok}}$ when the ciphertext
matches and $ss_{\mathrm{rej}}$ otherwise.  The function never returns
an error.  A protocol must not treat a returned 32-byte value as proof
that the ciphertext was valid; it must confirm the shared secret through
a subsequent authenticator or transcript binding.

\subsection{Correctness}

Ignoring compression, the decryption value is
\[
 c_2-\langle s,c_1\rangle
 = \Encode(m)+\langle e,r\rangle+e_2-\langle s,e_1\rangle.
\]
Compression adds residuals $\Delta_1$ on each $c_1$ coefficient and
$\Delta_2$ on $c_2$:
\[
 v'=\Encode(m)+\langle e,r\rangle+e_2-\langle s,e_1\rangle
 +\Delta_2-\langle s,\Delta_1\rangle.
\]
The residuals satisfy $\Delta_1\in\{-1,0,1,2\}$ and
$\Delta_2\in\{-63,\ldots,64\}$; over uniformly distributed coefficients
their variances are $1.25$ and $1365.25$.  A working variance model is:

\begin{table}[t]
\centering
\small
\begin{tabular}{@{}lr@{}}
\toprule
Component & Variance \\
\midrule
$\langle e,r\rangle$ & 1024 \\
$-\langle s,e_1\rangle$ & 1024 \\
$e_2$ & 1 \\
$-\langle s,\Delta_1\rangle$ & 1280 \\
$\Delta_2$ & 1365.25 \\
\midrule
Total & 4694.25 \\
Standard deviation & $\approx 68.5$ \\
\bottomrule
\end{tabular}
\caption{Working variance model for the KEM noise.  The independence and
Gaussian approximations are not proven.}
\label{tab:noise}
\end{table}

Rounding up gives $\sigma\approx 70$ and a margin of
$1024/70\approx 14.6\sigma$ against the decision boundary.  This is a
heuristic estimate, not a formal tail bound.  The decapsulation failure
probability has not been bounded to $2^{-128}$; the in-tree 32-trial
unit test is a regression tripwire, not the argument.  A formal
analysis must bound the tail of the bounded-sum distribution directly,
and a defensible empirical study must use at least $10^6$ trials with a
stated confidence interval.

\subsection{Intended security argument and its gaps}

The intended argument is FO-style: Module-LWE hides $m$ for a passive
adversary; the re-encryption check prevents ciphertext modification;
the random-oracle derivation of $(r,e_1,e_2)$ from $(m,pk)$ prevents
inconsistent randomness; implicit rejection removes a public validity
oracle.  The exact construction is not a verbatim copy of a standard FO
transform.  In particular, the accepted secret is
$\KDF(m\concat ct\concat\Hash(pk))$ and the rejection value is
$\KDF(\mathit{fake}\concat ct)$.  No cited theorem covers those exact
choices, and no reduction is proved in this repository.

The lattice parameters have not been run through the community
estimator.  The power-of-two modulus differs from ML-KEM-1024's 3329,
so ML-KEM cost estimates cannot be copied.  Consequently the KEM is
\emph{not} claimed IND-CCA2 secure; the target is conditional on the
missing analysis.

%======================================================================
\section{The CA-PRF and hash layer}
\label{sec:hash}

\subsection{CA-PRF construction}

For a 32-byte seed $x$, domain $d$, and output length $L$:
\[
 \begin{aligned}
 \mathit{init}&=\XOF(x\concat d\concat\texttt{"-init"},256),\\
 \mathit{rule}&=\XOF(x\concat d,64),\\
 \mathit{state}&=\mathrm{CA}^{32}(\mathit{init};\mathit{rule}),\\
 \mathit{out}&=\XOF(d\concat\mathit{state}\concat d\concat
 \texttt{"-out"},L).
 \end{aligned}
\]
The CA is a 32-word state with a fixed hyper-graph neighborhood
(strides 1, 5, 17), a secret 8-word rule, and 32 branchless rounds.
All loop bounds are public and all state is stack-resident.

The final BLAKE3 XOF is the security-relevant expander.  If BLAKE3 is a
secure XOF or random oracle, the output is pseudorandom regardless of
the CA.  The CA is therefore a design hedge and a deterministic mixing
stage, not an independently justified hardness assumption.  This
distinction is important: the paper does not claim that the automaton
adds provable security on top of BLAKE3.

\subsection{BLAKE3 API and framing}

The hash layer exposes unkeyed BLAKE3-256, a 64-byte XOF prefix, keyed
BLAKE3, the reference \texttt{derive\_key} construction, an expander
$\XOF(\mathit{seed}\concat\mathit{domain})$, a KDF
$\XOF(\mathit{domain}\concat\mathit{parts}\ldots)$, and a streaming
hasher.  \texttt{expand} and \texttt{kdf} do not insert length
separators; every protocol caller uses fixed-length fields or a single
variable-length field with a distinct domain.  The KEM and CA-PRF
callers satisfy this condition; a new caller must check it explicitly.

%======================================================================
\section{Authenticated encryption}
\label{sec:aead}

\subsection{One-shot construction}

Let $K$ be a 32-byte master key and $N$ a 12-byte nonce.  Derive
\[
 k_e=\mathrm{DK}(\texttt{"HYP-G/v07/symm-enc"},K),\qquad
 k_m=\mathrm{DK}(\texttt{"HYP-G/v07/symm-mac"},K).
\]
Encrypt $C=\mathrm{ChaCha20}(k_e,N,0,M)$ and authenticate
\[
 T=\mathrm{KeyedBLAKE3}(k_m,
 N\concat AD\concat C\concat\mathrm{LE64}(|AD|)\concat
 \mathrm{LE64}(|C|)).
\]
The one-shot opener recomputes the tag and performs a full-scan
constant-time comparison before writing any plaintext.  A tag mismatch
or length mismatch returns \texttt{None} and leaves the plaintext buffer
untouched.

The message limit is exactly $2^{38}$ bytes: $2^{32}$ ChaCha20 blocks of
64 bytes.  Counter exhaustion fails closed rather than wrapping.  The
nonce must be unique per $(K,\text{message})$; the construction has no
nonce-misuse resistance.

\subsection{Streaming front}

The streaming sealer buffers a partial 64-byte keystream block across
calls, feeds ciphertext bytes to a running keyed-BLAKE3 state, and
appends the two little-endian length suffixes at \texttt{finish}.  The
opener feeds ciphertext to the MAC while releasing tentative plaintext;
the plaintext becomes valid only when \texttt{finish} returns
\texttt{Some(())}.  On failure, the assembled plaintext buffer is
zeroized.  Empty chunks, one-byte chunks, block-boundary splits, and
1 MiB streams are tested.  The streaming output is byte-identical to
the one-shot output for the same inputs.

\subsection{Security position}

The EtM structure is standard, but the exact ChaCha20 + keyed-BLAKE3
composition is not a standardized AEAD and no reduction for this exact
instantiation is included.  The 32-byte tag is not truncated.  The
construction has no key-commitment guarantee; protocols requiring key
commitment must add it.  Nonce discipline is entirely the caller's
responsibility.

%======================================================================
\section{Hash-based signatures}
\label{sec:ds}

\subsection{Profile and keys}

The signature profile is $n=32$, $h=68$, $d=17$, $h'=4$, FORS
$(k,a)=(35,9)$, and WOTS+ $w=16$ with $\ell=67$.  The public key is
\[
 pk_{\mathrm{sig}}=\mathit{root}(32)\concat\mathit{pubseed}(32),
\]
and the secret key is
\[
 sk_{\mathrm{sig}}=\mathit{esk}(32)\concat\mathit{prf}(32)
 \concat\mathit{pubseed}(32)\concat\mathit{root}(32).
\]
The master secret is $\mathit{esk}$.  $\mathit{pubseed}$ is derived
from $\mathit{esk}$ with the \texttt{ds-pubseed} domain, $\mathit{prf}$
with \texttt{ds-prf}, and the root is the top XMSS layer root at layer
16, tree 0.

\subsection{Deterministic randomizer and digest}

The randomizer is
\[
 R=\Hash\bigl(\mathrm{CA\text{-}PRF}(\mathit{prf},
 \texttt{"HYP-G/v07/ds-addrnd"},64)\concat\Hash(\mathit{msg})\bigr).
\]
There is no fresh per-signature entropy.  Signing is a deterministic
function of $(sk,\mathit{msg})$, so repeated signing returns the same
bytes.  This is intentional for KATs and reproducibility, but it leaks
message equality and removes fresh-randomness fault hedging.  The scheme
is \emph{not} described as freshly hedged.

The 64-byte digest is
\[
 md_c=\Hash(R\concat\mathit{pubseed}\concat\mathit{root}
 \concat\Hash(\mathit{msg})\concat\mathrm{LE32}(c)
 \concat\texttt{"HYP-G/v07/ds-msghash"})
\]
for $c\in\{0,1\}$.  The first 387 bits are split into 35 nine-bit FORS
indices, a 68-bit tree index, and a four-bit bottom leaf.  The remaining
bits are unused.

\subsection{Tweakable hash and v2.0.1 tags}

All tweakable hashes are keyed BLAKE3 with $\mathit{pubseed}$:
\[
 \begin{aligned}
 \PRF(\mathit{pubseed},x,A)&=\mathrm{KeyedBLAKE3}
   (\mathit{pubseed},01\concat x\concat A),\\
 F(\mathit{pubseed},A,x)&=\mathrm{KeyedBLAKE3}
   (\mathit{pubseed},02\concat A\concat x),\\
 \Hash(\mathit{pubseed},A,L,R)&=\mathrm{KeyedBLAKE3}
   (\mathit{pubseed},A\concat L\concat R).
 \end{aligned}
\]
The explicit tags are the DS 2.0.1 revision.  Before it, $\PRF$ and
$F$ both had 64-byte inputs and could be structurally identical when a
secret seed happened to equal the address bytes.  The tags make the
first input byte different in every call.  The 96-byte untagged $H$
call has a distinct arity and field order.

\subsection{Trees}

FORS uses 35 trees of height 9.  Each tree signs one nine-bit leaf
index; the signature contains the selected secret leaf and its nine
authentication nodes (320 bytes per tree).  WOTS+ maps a 32-byte message
to 64 nibbles plus a three-digit checksum, and signs by applying a
15-step chain; every chain executes all 15 steps with a branchless
select so the secret digit does not control the loop trip count.  The
L-tree compresses 67 chain ends with seven levels; an unpaired final
node is promoted unchanged.  An XMSS layer is a height-4 Merkle tree
over 16 WOTS+ leaves.  The hypertree stacks 17 such layers; tree
$i$ is $\mathit{tree\_idx}\gg 4(i+1)$, the bottom leaf is the digest's
four-bit leaf, and upper leaf indices are the low four bits of the
child tree index.

\subsection{Encoding and verification}

The 49,992-byte signature is
\[
 R(32)\concat \mathrm{FORS}_{35}\concat
 \mathrm{layer}_0\concat\cdots\concat\mathrm{layer}_{16}.
\]
Each layer is an eight-byte reserved zero field, a 67-element WOTS+
signature, and a four-node authentication path.  Verification checks the
exact length, rejects any nonzero reserved byte, recomputes the digest
and FORS key, climbs all layers, and compares the resulting root with
the public root using constant-time equality.  No field is silently
ignored.

\subsection{Security position}

The composition is SLH-DSA-shaped but not FIPS 205 conformant.  It has
no exact UF-CMA/SUF-CMA proof for its BLAKE3 instantiation, address
layout, deterministic randomizer, or reserved-byte encoding.  FORS
is a few-time signature: security degrades as more signatures are
produced under one key.  The implementation enforces no signing limit
and keeps no state.  A deployment profile would need a concrete bound on
signatures per key and a matching multi-target analysis.  The current
paper claims only that the implementation computes the specified bytes
and that the deterministic vectors and tamper tests pass.

%======================================================================
\section{Implementation assurance}
\label{sec:implementation}

\subsection{Memory and allocation discipline}

The cryptographic libraries use fixed-size arrays and caller-provided
buffers.  They do not allocate, do not read the clock, do not access the
filesystem, and do not spawn threads.  The \texttt{std} crates own all
OS interaction.  CI greps the \texttt{no\_std} libraries for
\texttt{Vec}, \texttt{Box}, \texttt{alloc::}, \texttt{String}, and
\texttt{std::} and fails if any appears.

\subsection{Secret lifetimes}

Secret-holding types redact \texttt{Debug}, avoid \texttt{Copy} where
duplication would be dangerous, compare in constant time where equality
is exposed, and best-effort zeroize on \texttt{Drop}.  The types include
the KEM secret key and shared secret, the DS secret key and pre-sign
context, the AEAD key, the CA state and rule, the hash state, and
polynomial buffers.  This is best-effort erasure in safe Rust: normal
writes are followed by compiler fences and a black-box use of the
buffer pointer, but safe Rust cannot guarantee volatile erasure.  The
accessors \texttt{seed\_copy}, \texttt{fake\_copy}, and
\texttt{public\_copy} return caller-owned copies that the caller must
erase.

\subsection{Constant-time posture}

The implementation avoids secret-dependent branches and indices in the
cryptographic paths: KEM decapsulation always performs re-encryption and
both KDFs; WOTS+ chains always execute all 15 steps; tag and root
comparisons scan all bytes; implicit-rejection selection is a byte mask.
This is a source-level discipline, not a machine-code guarantee.
Compiler transformations and microarchitectural effects are not
 analysed, and no \texttt{dudect}~\cite{dudect} or power/EM/fault
campaign has been performed.

\subsection{Failure semantics}

The KEM never returns an error; invalid ciphertexts yield a
pseudorandom secret.  The one-shot AEAD returns \texttt{None} and leaves
the output untouched on authentication failure.  The streaming opener
marks plaintext as tentative until \texttt{finish}.  Signature
verification returns \texttt{false} on any malformation.  The CLI uses
exit code 2 for usage/input errors and 3 for authentication failures.

%======================================================================
\section{Evaluation}
\label{sec:evaluation}

\subsection{Method and evidence levels}

We use five evidence levels: E0 source inspection; E1 deterministic
tests and KATs; E2 adversarial harnesses; E3 dynamic/portability
analysis such as Miri and cross-architecture runs; E4 independent
review.  The current artifact has E0--E3 evidence and no E4 evidence.
No result below is an external audit.

The controlled benchmark command was:
\begin{verbatim}
HYPG_BENCH_ITERS=8 \
HYPG_BENCH_DS_ITERS=8 \
HYPG_BENCH_TTEST_N=10000 \
cargo run -q --release -p hypg-bench
\end{verbatim}
The machine was an Intel Core i5-14400F with 16 hardware threads and
31 GiB RAM, running Linux x86-64 with \texttt{rustc 1.97.1}.  The
benchmark warms up once and reports integer averages over eight
iterations; large-message rates use three repetitions of 1 MiB.

\subsection{Correctness and reproducibility}

The debug, release, and MSRV 1.70 test suites each report
\textbf{146 passed, 0 failed, 0 ignored}.  The breakdown is 108 unit
tests, 25 CLI negative integration tests, and 13 doc-tests.  The
red-team harness passes all 14 suites.  The frozen KAT inventory is:

\begin{table}[t]
\centering
\small
\begin{tabular}{@{}lrr@{}}
\toprule
File & Bytes & SHA-256 (prefix) \\
\midrule
\texttt{kats/v2.0.0.txt} & 28,514 & \texttt{3b2dde3d...71d4b} \\
\texttt{kats-symm/v2.0.0.txt} & 3,662 & \texttt{91b7942f...4208} \\
\texttt{kats-hash/v2.0.0.txt} & 923 & \texttt{5fa77803...d9f1} \\
\texttt{kats-ds/v2.0.1.txt} & 201,015 & \texttt{e88c1bdf...c813} \\
\midrule
Total & 234,114 & \\
\bottomrule
\end{tabular}
\caption{Frozen vector inventory.  The complete hashes are in
\texttt{docs/parameters.md} and \texttt{docs/testplan.md}.}
\label{tab:kats}
\end{table}

The complete SHA-256 values are:

\begin{verbatim}
kats/v2.0.0.txt
3b2dde3d1138c11b54a6880f50ca7b4390071ae53aa94601075f39b2af971d4b
kats-symm/v2.0.0.txt
91b7942ffb4a3c5da8b8b589268af7b7d1594796b428efdc7e26c567c5f04208
kats-hash/v2.0.0.txt
5fa778035ce2e8aa9c89345458ba4da67df4bdbe71a456cbb58026cf9341d9f1
kats-ds/v2.0.1.txt
e88c1bdf54e9390736cd03c1a165616d744a42c48d662880b344c7de3546c813
\end{verbatim}

The x86-64 and \texttt{qemu-aarch64} static-musl KAT outputs are
byte-identical across all 234,114 bytes.  This checks endianness,
integer width, and the DS 2.0.1 revision on an emulated ARM64 target;
it is not native ARM64 validation.

\subsection{Adversarial testing}

The deterministic red-team harness covers garbage inputs, ciphertext
malleability, an $8\times8$ key-mismatch matrix, message extremes, FO
binding, sampler statistics over 524,288 CBD coefficients, failure
sweeps, determinism, AEAD forgery, DS forgery, burst faults, tampered
public keys, padding flips, a panic sweep, secret-key bit flips, and
streaming AEAD games.  A fresh release run with 20,000 trials reported
20,000/20,000 clean decapsulations.  A separate full-bit run flipped all
11,520 bit positions of a test ciphertext; every flip changed the
shared secret.  These are regression and sanity checks, not a
$2^{-128}$ failure proof.

\subsection{Miri and portability}

Miri is an interpreter for Rust's mid-level IR.  Latest recorded runs
on the hardened commit report KEM 9/9 (13,213.79 s), AEAD 6/6
(214.46 s), and DS bounds 3/3 (0.42 s).  Full DS signing is not run
under Miri because the hypertree is hours-long under interpretation.
Miri checks executed paths for some classes of undefined behaviour; it
does not prove their absence.  The libraries also build for
\texttt{aarch64}, \texttt{wasm32}, \texttt{s390x} (check only), and
\texttt{thumbv7em}; KAT and red-team binaries pass under a 1 MB stack.

\subsection{Performance}

\begin{table}[t]
\centering
\small
\begin{tabular}{@{}lr@{}}
\toprule
Operation & Controlled result \\
\midrule
KEM keygen / encaps / decaps & 795 / 931 / 1124 $\mu$s \\
BLAKE3 hash / streaming / XOF (1 MiB) & 842 / 837 / 827 MB/s \\
ChaCha20 XOR (1 MiB) & 773 MB/s \\
AEAD seal / open (1 MiB) & 427 / 434 MB/s \\
DS sign / verify & 72.009 ms / 3.725 ms \\
Threaded DS sign (16 workers) & 18.329 ms, 3.928$\times$ \\
\bottomrule
\end{tabular}
\caption{Controlled release measurements on the record machine.
These are local point estimates, not portable guarantees.}
\label{tab:bench}
\end{table}

The valid-versus-tampered medians were 1004 per-mille for KEM
decapsulation, 602 per-mille for AEAD open, and 994 per-mille for DS
verify.  The Welch t-test over 10,000 interleaved KEM samples had
$\max |t|=0.737$.  These are timing-smoke signals, not a side-channel
proof.

\subsection{Interpretation}

The evaluation establishes that the specified bytes are reproducible,
that the tested failure paths behave as designed, that the KATs survive
debug/release and cross-architecture comparison, and that the scalar
implementation has usable research performance.  It does not establish
cryptographic security, constant-time execution, or production
readiness.

%======================================================================
\section{Security analysis and open problems}
\label{sec:security}

\subsection{Assumptions}

The intended security argument depends on: (A1) Module-LWE hardness for
the exact parameters; (A2) BLAKE3 as a collision-resistant hash,
PRF/MAC, and XOF; (A3) ChaCha20 as a secure stream cipher with unique
nonces; (A4) the exact FO-style transform having CCA security under
A1--A2; (A5) random-oracle behaviour of the hash/KDF calls; (A6) a
negligible decapsulation-failure probability; (A7) nonce uniqueness;
(A8) a bounded DS signing volume; (A9) caller-protected secret storage;
and (A10) no machine-code constant-time claim.

\subsection{The exact FO variant}

Standard FO analyses~\cite{fo99,hhk17} do not automatically cover the
chosen shared-secret and rejection derivations.  A proof must adapt the
transform, account for the full public-key binding in $G$, and bound
the re-encryption-collision probability.  This is an open item, not a
footnote.

\subsection{The DS signing volume}

FORS is few-time.  The security of the hypertree depends on how many
messages are signed under one key and on the probability of reusing
FORS leaves or WOTS+ keys.  The current implementation stores no state
and enforces no limit.  A deployment profile needs a concrete bound and
a multi-target analysis; neither exists here.

\subsection{The AEAD composition}

The EtM composition is structurally standard, but the exact
ChaCha20 + keyed-BLAKE3 instantiation has no cited reduction.  Nonce
reuse is catastrophic and caller-controlled.  No key-commitment property
is claimed.

\subsection{Side channels and fault injection}

Source-level branchlessness is not a machine-code guarantee.  The
project has no \texttt{dudect} evaluation, no power/EM laboratory, and
no fault-injection campaign.  The deterministic DS randomizer also
removes fresh-randomness fault hedging.  These are material limitations
for any deployment-oriented claim.

\subsection{Operational assumptions}

The caller must supply a CSPRNG, protect keys in memory, ensure nonce
uniqueness, confirm KEM secrets, bound signature volume, and provide
protocol framing and replay protection.  The libraries provide none of
these services.

%======================================================================
\section{Threats to validity}
\label{sec:validity}

\paragraph{Construct validity.}  KATs and reproducibility checks measure
\emph{behaviour}, not hardness.  A deterministic test can pass while the
construction is insecure.

\paragraph{Internal validity.}  The red-team harness is deterministic and
finite.  It does not cover every input, every protocol composition, or
every compiler transformation.  The 20,000-trial failure sweep is not a
tail bound.

\paragraph{External validity.}  Performance is measured on one CPU,
one toolchain, and one operating system.  Cross-architecture KATs are
emulated on ARM64; native ARM64 and s390x runtime validation are absent.
The results do not establish portable timing behaviour.

\paragraph{Conclusion validity.}  The absence of external audit,
lattice-estimator output, formal proof, and side-channel data prevents
any conclusion about deployable security.

%======================================================================
\section{Related work}
\label{sec:related}

ML-KEM~\cite{fips203}, ML-DSA~\cite{fips204}, and
SLH-DSA~\cite{fips205} are the NIST standards closest to the relevant
families.  HypG shares the Module-LWE and hypertree shapes but is not
wire-compatible or standards-conformant.  SPHINCS+~\cite{sphincs} and
XMSS~\cite{rfc8391} are the main references for the signature
composition.  BLAKE3~\cite{blake3}, ChaCha20~\cite{rfc8439}, and the
Fujisaki--Okamoto literature~\cite{fo99,hhk17} are the symmetric and
transform references.  The lattice-estimator~\cite{estimator} is the
required external tool for the KEM parameters.  The
\texttt{dudect}~\cite{dudect} methodology is the minimum bar for a
timing evaluation.  Rust's Miri~\cite{miri} is used for partial
undefined-behaviour checking.

HypG's contribution is not a new hardness assumption.  It is an
integrated, transparent, reproducible research artifact with an
explicit evidence boundary.

%======================================================================
\section{Road to a defensible profile}
\label{sec:roadmap}

The shortest honest path to a deployable profile is:

\begin{enumerate}[leftmargin=*]
\item Run the community lattice-estimator on the exact
  $(n,k,q,\eta)=(256,4,4096,2)$ parameters, publish classical and
  quantum cost tables, and compare against ML-KEM-1024 without claiming
  equivalence.
\item Derive an analytic decapsulation-failure bound and measure it at
  at least $10^6$ trials.
\item Prove or refute the exact FO-style transform, including the
  accepted and rejected secret derivations.
\item Prove a UF-CMA/SUF-CMA bound for the exact DS composition and
  choose an enforced signing-volume limit.
\item Complete a \texttt{dudect}-style timing evaluation and, if the
  threat model requires it, power/EM/fault evaluation.
\item Add coverage-guided fuzzing and a protocol-level misuse analysis.
\item Commission an independent code audit and publish findings.
\item If FIPS compliance is required, redesign around approved
  primitives and pursue ACVP/CAVP and CMVP validation.
\end{enumerate}

Until then, HypG should remain in research, teaching, reproducibility,
and test environments.

%======================================================================
\section{Conclusion}
\label{sec:conclusion}

HypG v2.0 is a complete, zero-external-dependency, safe-Rust
post-quantum research suite with a documented KEM, hash layer, AEAD, and
stateless signature scheme.  Its evaluation is unusually broad for a
prototype: frozen KATs, cross-architecture equality, 146 tests, a
14-suite adversarial harness, Miri subsets, and controlled performance
measurements.  Its security claims are deliberately narrow: the
artifact implements the specification, the tested paths behave as
specified, and the remaining proof, estimator, side-channel, and audit
work is named explicitly.

That combination---substantial engineering evidence and an explicit
refusal to round it up into a security claim---is the contribution of
this paper.

\appendix

%======================================================================
\section{Parameter register}
\label{app:params}

\begin{longtable}{@{}lll@{}}
\toprule
Object & Value & Note \\
\midrule
\endhead
KEM ring & $n=256$, $q=4096$ & $R_q=\Z_q[X]/(X^{256}+1)$ \\
KEM rank & $k=4$ & lattice dimension 1024 \\
KEM noise & $\mathrm{CBD}(2)$ & variance 1, support $\{-2,\ldots,2\}$ \\
KEM sizes & $pk=1568$, $ct=1440$, $sk=1632$, $ss=32$ & bytes \\
DS hash & $n=32$ & BLAKE3-256 \\
DS tree & $h=68$, $d=17$, $h'=4$ & 17 XMSS layers \\
DS FORS & $k=35$, $a=9$ & 35 trees, 512 leaves each \\
DS WOTS+ & $w=16$, $\ell=67$ & 64 message + 3 checksum digits \\
DS sizes & $pk=64$, $sk=128$, $sig=49992$ & bytes \\
AEAD & key 32, nonce 12, tag 32 & bytes \\
AEAD limit & $2^{38}$ & bytes per $(key,nonce)$ \\
CA state & 32 words of 64 bits & 256 bytes \\
CA rounds & 32 & fixed \\
\bottomrule
\end{longtable}

%======================================================================
\section{Artifact and reproducibility notes}
\label{app:artifact}

The active documentation is in \texttt{docs/}.  The frozen vectors are
under \texttt{crates/hypg-kat/}.  The controlled benchmark command and
machine details are in \texttt{docs/metrics.md}.  The test matrix and
evidence levels are in \texttt{docs/testplan.md}.  The exact security
non-claims are in \texttt{docs/threat-model.md} and
\texttt{SECURITY.md}.

The core reproducibility commands are:

\begin{verbatim}
cargo test --workspace --locked
cargo test --workspace --locked --release
cargo +1.70 test --workspace --locked
cargo run -q -p hypg-redteam
cargo run -q -p hypg-lattice
cargo test -p hypg-kat --lib --locked
HYPG_BENCH_ITERS=8 HYPG_BENCH_DS_ITERS=8 HYPG_BENCH_TTEST_N=10000 \
  cargo run -q --release -p hypg-bench
\end{verbatim}

The code is licensed under GPL-3.0-or-later.  The paper sources are
tracked in \texttt{papers/}; the PDFs are built with \texttt{pdflatex}.

\begin{thebibliography}{99}

\bibitem{fips203}
National Institute of Standards and Technology.
\newblock \emph{Module-Lattice-Based Key-Encapsulation Mechanism
Standard}. FIPS 203, August 2024.

\bibitem{fips204}
National Institute of Standards and Technology.
\newblock \emph{Module-Lattice-Based Digital Signature Standard}.
FIPS 204, August 2024.

\bibitem{fips205}
National Institute of Standards and Technology.
\newblock \emph{Stateless Hash-Based Digital Signature Standard}.
FIPS 205, August 2024.

\bibitem{blake3}
J.~O'Connor, J.-P. Aumasson, S.~Neves, and Z.~Wilcox-O'Hearn.
\newblock BLAKE3: one function, fast everywhere.
\newblock Technical report, 2020.
\newblock \url{https://github.com/BLAKE3-team/BLAKE3-specs/}.

\bibitem{rfc8439}
Y.~Nir and A.~Langley.
\newblock ChaCha20 and Poly1305 for IETF Protocols.
\newblock RFC 8439, June 2018.

\bibitem{fo99}
E.~Fujisaki and T.~Okamoto.
\newblock Secure integration of asymmetric and symmetric encryption
schemes.
\newblock In \emph{CRYPTO 1999}, LNCS 1666, pages 537--554, 1999.

\bibitem{hhk17}
D.~Hofheinz, K.~H\"ovelmanns, and E.~Kiltz.
\newblock A modular analysis of the Fujisaki--Okamoto transformation.
\newblock In \emph{TCC 2017}, LNCS 10677, pages 341--371, 2017.

\bibitem{sphincs}
D.~J. Bernstein, A.~H\"ulsing, S.~K\"olbl, R.~Niederhagen,
J.~Rijneveld, and P.~Schwabe.
\newblock The SPHINCS+ signature framework.
\newblock In \emph{ACM CCS 2019}, pages 2129--2146, 2019.

\bibitem{rfc8391}
A.~H\"ulsing, D.~Butin, S.-L. Gazdag, J.~Rijneveld, and A.~Mohaisen.
\newblock XMSS: eXtended Merkle Signature Scheme.
\newblock RFC 8391, May 2018.

\bibitem{estimator}
M.~R. Albrecht, R.~Player, and S.~Scott.
\newblock On the concrete hardness of Learning with Errors.
\newblock \emph{Journal of Mathematical Cryptology}, 9(3):169--203,
2015.

\bibitem{dudect}
O.~Reparaz, J.~Balasch, and I.~Verbauwhede.
\newblock Dude, is my code constant time?
\newblock In \emph{DATE 2017}, pages 1697--1702, 2017.

\bibitem{miri}
The Rust Project.
\newblock Miri: an interpreter for Rust's mid-level intermediate
representation.
\newblock \url{https://github.com/rust-lang/miri}.

\end{thebibliography}

\end{document}
